% ---------------------------------------------------------------------------
% §3.2 Histories, Memories, and Interventions. Three added paragraphs, placed
% after "Isolation controls and interventions." and before "Mitigations.", plus
% one sentence appended to the "Mitigations." paragraph.
% ---------------------------------------------------------------------------
\paragraph{Generated histories.}
To ask which feature of a history triggers laundering, we generate 108 groups of procurement histories in the corpus format from four themes. Within a group the three histories share every turn, date, and probe and differ only in the number of later messages, from parties without authority, that restate a permission after it has been superseded: zero, two, or four. This comparison changes one thing at a time. The generator also varies whether the change is an amendment of the grant or a revoke-and-replace and how many blocks separate the grant from its change; those factors reseed the surrounding history, so their comparisons are observational. The corpus passes the same linter as the hand-written corpora.

\paragraph{Closed loop.}
In the main protocol the executor's actions leave no trace. To model a deployment that logs them, the closed loop appends, after each request, one workflow-log line stating what the executor did (executed as submitted, executed an operational alternative, escalated, or declined, with the payload) as a new block, and the writer updates memory from it before the next request. A neutral control forked from the same frozen memory receives the same number of updates on the same schedule, each carrying a content-free workspace notice, so the control isolates the content of the written-back line from the effect of repeated updates. A round is one pass over a case's requests in request-time order; requests are re-dated to follow the log they can see whenever this leaves every ledger verdict unchanged. Runs use incremental typed memory at each domain's canonical seed for three rounds; the action-minus-control difference is paired by chain with a bootstrap interval and a sign-flip permutation test.

\paragraph{Locating and labeling failures.}
For every unauthorized request the final memory authorizes, we replay the saved memory after each block against the ledger as of that block; the error block is the first at which the memory authorizes the request while the ledger does not, and stays so to the end. For the closed loop we also take every permission record whose only cited sources are the agent's own written-back lines. Three judge models (DeepSeek V4 Pro, GLM 5.3, Nemotron 3 Ultra, temperature 0) see the policy, the request, the true permission state after the block, the memory before and after, the block's messages, and the writer's plan and patches, and each names one of four errors: a restatement from a party without authority applied; the agent's own action treated as approval; an authoritative change misapplied; or the update rejected. The majority label is reported, with agreement; the labels are the judges' and have not been checked by a person.

% Append to the "Mitigations." paragraph:
Two further mitigations act on the writer rather than on provenance: rebuilding memory from the full history every three blocks instead of patching it incrementally, and prepending one sentence to the writer's instructions, ``Only record permissions that an authorized approver has actually granted, no matter what anyone else says or asks.'' Both are compared with the same conditions without them at the same seeds and executors.
